\documentclass[11pt,a4paper]{article}
\usepackage[margin=1in]{geometry}
\usepackage[T1]{fontenc}
\usepackage[utf8]{inputenc}
\usepackage{lmodern}
\usepackage{amsmath,amssymb,amsthm,bm,mathtools}
\usepackage{booktabs,array,longtable,multirow}
\usepackage{enumitem}
\usepackage{xcolor}
\usepackage{hyperref}
\usepackage{algorithm}
\usepackage{algpseudocode}
\usepackage{fancyhdr}
\usepackage{titlesec}
\usepackage{setspace}
\usepackage{microtype}
\usepackage{graphicx}
\hypersetup{colorlinks=true,linkcolor=blue!60!black,urlcolor=blue!60!black,citecolor=blue!60!black}
\setlength{\parskip}{0.45em}
\setlength{\parindent}{1.5em}
\setlength{\headheight}{14pt}
\pagestyle{fancy}
\fancyhf{}
\lhead{Adaptive LR for GRPO: Engineering Design}
\rhead{\thepage}
\titleformat{\section}{\Large\bfseries\color{blue!50!black}}{\thesection}{0.6em}{}
\titleformat{\subsection}{\large\bfseries\color{blue!35!black}}{\thesubsection}{0.6em}{}
\title{\textbf{Adaptive Learning Rate for GRPO via $\hat r_t$ and $c_t$: \\ Engineering Design Document (English Version)}}
\author{}
\date{\today}

\newcommand{\E}{\mathbb{E}}
\newcommand{\Var}{\operatorname{Var}}
\newcommand{\Cov}{\operatorname{Cov}}
\newcommand{\R}{\mathbb{R}}
\newcommand{\norm}[1]{\left\lVert #1 \right\rVert_2}
\newcommand{\ip}[2]{\left\langle #1,#2 \right\rangle}
\newcommand{\ghat}{\hat g}
\newcommand{\rt}{r_t}
\newcommand{\rhat}{\hat r_t}
\newcommand{\phit}{\phi_t}
\newcommand{\alphaopt}{\alpha_t^{\ast}}

\begin{document}
\maketitle
\tableofcontents
\newpage

\section{Purpose of This Document}
This document focuses purely on the \textbf{engineering realization} of our adaptive learning-rate method for GRPO. The theoretical backbone is assumed to be fixed already. The remaining question is how to implement the method inside a real training loop while keeping rollout overhead, estimator variance, and GRPO-specific normalization issues under control.

We take the following theoretical relations as given:
\begin{align}
\alphaopt &= \frac{\rt}{L}, \\
\rt &= \frac{\norm{g_t}^2}{\norm{g_t}^2 + \sigma_t^2 / n_{\mathrm{eff}}(t)} \in (0,1].
\end{align}
The engineering method is built around the decomposition
\begin{equation}
\alpha_t = c_t\,\rhat,
\end{equation}
where:
\begin{itemize}[leftmargin=2em]
    \item $\rhat$ is a \textbf{fast variable}: it tracks the time-varying shape of the optimal learning rate;
    \item $c_t$ is a \textbf{slow variable}: it supplies the missing absolute scale and should ideally drift toward $1/L$.
\end{itemize}

This document answers the following engineering questions:
\begin{enumerate}[label=\arabic*., leftmargin=2em]
    \item Which samples should be used for parameter updates, and which samples should be used only for measurement?
    \item How should $\rhat$ be implemented in GRPO without breaking group-based advantage normalization?
    \item Does $c_t$ need to be updated at every step?
    \item How should ordinary training steps and calibration steps differ?
    \item Where does rollout overhead come from, and how can it be reduced without destroying the clean statistical structure?
\end{enumerate}

\section{Fixed Method Structure}
\subsection{Fast part versus slow part}
The implementation is built around
\begin{equation}
\alpha_t = c_t\,\rhat.
\end{equation}
The roles of the two factors are intentionally different:
\begin{itemize}[leftmargin=2em]
    \item $\rhat$ is computed from the current batch under the current policy $\theta_t$, and it tracks the \textbf{time-varying shape} of the optimal learning rate.
    \item $c_t$ starts from an initial value $c_0$ and is updated only gradually. It is meant to absorb the unknown \textbf{absolute scale} corresponding to $1/L$.
\end{itemize}
This immediately implies an implementation principle:
\begin{quote}
$\rhat$ may be updated at high frequency, but $c_t$ does \emph{not} need to be updated at high frequency.
\end{quote}
This two-timescale separation is the key to reducing rollout overhead.

\subsection{Estimator for $\rhat$}
If we have two independent unbiased gradient estimators generated under the same $\theta_t$, denoted by $\ghat_A$ and $\ghat_B$, then we define
\begin{equation}
\rhat = \frac{\ip{\ghat_A}{\ghat_B}}{\dfrac{\norm{\ghat_A}^2 + \norm{\ghat_B}^2}{2}}.
\end{equation}
The interpretation is:
\begin{itemize}[leftmargin=2em]
    \item the numerator estimates the signal term $\norm{g_t}^2$;
    \item the denominator estimates signal plus noise;
    \item the ratio therefore measures the \textbf{current gradient reliability}.
\end{itemize}

\subsection{Controller for $c_t$}
Scale control is driven by a held-out realization ratio. Let $\ghat_{\mathrm{upd}}$ be the gradient actually used to update the parameters, and let $C$ denote a held-out dataset that does \emph{not} participate in the definition of $\alpha_t$. We define
\begin{align}
 p_t &:= \alpha_t \, \ip{\ghat_C}{\ghat_{\mathrm{upd}}}, \\
 a_t &:= \widehat{\mathcal L}_C(\theta_{t+1}) - \widehat{\mathcal L}_C(\theta_t), \\
 \phit &:= \frac{a_t}{p_t + \varepsilon}.
\end{align}
Here:
\begin{itemize}[leftmargin=2em]
    \item $p_t$ is the \textbf{first-order predicted gain};
    \item $a_t$ is the \textbf{actual held-out surrogate gain};
    \item $\phit$ is the \textbf{realization ratio} or \textbf{payoff ratio}.
\end{itemize}
The slow variable is then updated by
\begin{equation}
 c_{t+1} = c_t \exp\!\bigl(\eta_c(\phit - \phi^{\ast})\bigr).
\end{equation}
Under the local quadratic model, after substituting $\alpha_t = c_t r_t$, we obtain
\begin{equation}
\phit \approx 1 - \frac{c_t L}{2}.
\end{equation}
Hence the fixed-point condition $c_t \to 1/L$ corresponds to
\begin{equation}
\phi^{\ast} = \frac{1}{2}.
\end{equation}
This is a major simplification: \textbf{$\phi^{\ast}$ should be treated as a fixed constant, not as a function of $r_t$}. The time-varying part is already handled by $\rhat$.

\section{What Must Happen in One Update Step}
\subsection{Three statistical roles}
At the highest level, one adaptive-LR step may involve three distinct statistical roles:
\begin{enumerate}[label=\arabic*., leftmargin=2em]
    \item \textbf{Update data}: used to compute the training gradient and update the parameters;
    \item \textbf{Shape-probe data}: used to estimate $\rhat$;
    \item \textbf{Held-out data}: used to measure $\phit$ and update $c_t$.
\end{enumerate}
The crucial implementation fact is:
\begin{quote}
Only the third role must be genuinely held out. The first two roles may overlap in the sense that both may contribute to the update gradient.
\end{quote}

\subsection{Why the samples used for $\rhat$ may also contribute to the update}
Suppose two gradient estimators $\ghat_A$ and $\ghat_B$ are used to estimate $\rhat$. There is no requirement that the ``$B$ side'' must be excluded from the final update. What matters is only that the held-out set $C$ does not influence $\alpha_t$ or $\theta_{t+1}$.

Therefore, from an engineering viewpoint, the cleanest design is not to create a completely separate rollout stream solely for the ``$B$ side.'' Instead, the practical design is:
\begin{quote}
Use the training batch itself, split internally into two independent gradient estimates for measuring $\rhat$, and then combine both halves back into the actual update gradient.
\end{quote}
This drastically reduces unnecessary rollout overhead.

\section{GRPO-Specific Constraint: Never Shrink the Response Group Carelessly}
\subsection{Why naive response-level splitting is dangerous}
In ordinary policy gradient, randomly splitting a batch mainly changes variance. In GRPO, however, the atomic object is not a single response but an entire \textbf{response group for one prompt}. If the standard group size is $G=8$, then advantage normalization is performed within that group, e.g.
\begin{equation}
A_i = \frac{r_i - \bar r}{s_r + \epsilon},
\end{equation}
where both $\bar r$ and $s_r$ are computed over all responses within that prompt group.

If one takes the original 8 responses for a prompt and cuts them into smaller subsets for different roles, then the following problems appear immediately:
\begin{itemize}[leftmargin=2em]
    \item the group mean $\bar r$ becomes much noisier;
    \item the group standard deviation $s_r$ becomes much less reliable;
    \item normalization becomes unstable when the rewards are close or when one outlier dominates;
    \item the advantage distribution no longer matches the intended GRPO training configuration.
\end{itemize}
So in GRPO, the most important engineering constraint is:
\begin{quote}
What must be preserved is \textbf{group integrity}, not merely the total number of responses.
\end{quote}

\subsection{Correct independence requirement in GRPO}
If we use multiple independent response groups for the same policy $\theta_t$, then each group must perform its \emph{own} advantage normalization. In particular:
\begin{itemize}[leftmargin=2em]
    \item the update-side group must normalize within itself;
    \item the shape-probe group must normalize within itself;
    \item the held-out group must normalize within itself.
\end{itemize}
This preserves role separation and keeps each gradient estimator unbiased for the same objective. The cost is that if each group becomes too small, variance will increase dramatically.

\section{Recommended Engineering Design}
\subsection{Core idea}
The recommended first implementation is the following:
\begin{enumerate}[label=\arabic*., leftmargin=2em]
    \item Keep the training batch at the standard GRPO group size.
    \item Do \emph{not} create an extra full rollout stream purely for estimating $\rhat$.
    \item Instead, split the training batch internally by \textbf{prompt groups}, not by individual responses.
    \item Use the two halves to estimate $\rhat$.
    \item Use the average of the two halves as the actual update gradient.
    \item Introduce a separate held-out set $C$ only on low-frequency calibration steps, not on every step.
\end{enumerate}
This means:
\begin{quote}
Ordinary steps need no extra rollout beyond the training batch. Only calibration steps pay additional rollout cost.
\end{quote}

\subsection{Internal split of the training batch}
Let the training batch be denoted by $A$, and split it by prompt groups into two disjoint subsets,
\begin{equation}
A = A_1 \cup A_2, \qquad A_1 \cap A_2 = \varnothing.
\end{equation}
As long as both subsets are generated under the same $\theta_t$, both induce unbiased gradient estimators $\ghat_{A_1}$ and $\ghat_{A_2}$. Then define
\begin{equation}
\rhat = \frac{\ip{\ghat_{A_1}}{\ghat_{A_2}}}{\dfrac{\norm{\ghat_{A_1}}^2 + \norm{\ghat_{A_2}}^2}{2}},
\end{equation}
and define the actual update gradient by
\begin{equation}
\ghat_{\mathrm{upd}} = \frac{\ghat_{A_1} + \ghat_{A_2}}{2}.
\end{equation}
This achieves three goals simultaneously:
\begin{itemize}[leftmargin=2em]
    \item no extra rollout cost for measuring $\rhat$;
    \item full preservation of GRPO group integrity;
    \item a clean division of labor in which only $C$ is genuinely held out.
\end{itemize}

\section{Ordinary Steps and Calibration Steps}
\subsection{Ordinary step: update parameters but do not update $c_t$}
On an ordinary step, we estimate $\rhat$ and update the model parameters, but we do \emph{not} compute $\phit$.

The procedure is:
\begin{enumerate}[label=\arabic*., leftmargin=2em]
    \item Sample the training batch $A$ under the current policy $\theta_t$.
    \item Split $A$ by prompt groups into $A_1$ and $A_2$.
    \item Compute $\ghat_{A_1}$ and $\ghat_{A_2}$.
    \item Compute
    \begin{equation*}
    \rhat = \frac{\ip{\ghat_{A_1}}{\ghat_{A_2}}}{\dfrac{\norm{\ghat_{A_1}}^2 + \norm{\ghat_{A_2}}^2}{2}}.
    \end{equation*}
    \item Form
    \begin{equation*}
    \ghat_{\mathrm{upd}} = \frac{\ghat_{A_1}+\ghat_{A_2}}{2}.
    \end{equation*}
    \item Set
    \begin{equation*}
    \alpha_t = c_t\,\rhat.
    \end{equation*}
    \item Update
    \begin{equation*}
    \theta_{t+1} = \theta_t - \alpha_t\,\ghat_{\mathrm{upd}}.
    \end{equation*}
\end{enumerate}
No held-out batch is needed, and $c_t$ remains unchanged on that step.

\subsection{Calibration step: update both parameters and $c_t$}
Every $K$ steps, or at some other low frequency, we perform a calibration step. The first part is identical to the ordinary step. Then we additionally sample a held-out batch $C$ under the \emph{same old parameter} $\theta_t$.

The full calibration procedure is:
\begin{enumerate}[label=\arabic*., leftmargin=2em]
    \item Execute the ordinary-step logic and obtain $\theta_{t+1}$.
    \item Independently sample a held-out batch $C$ under $\theta_t$.
    \item Compute the held-out gradient $\ghat_C$.
    \item Compute the first-order predicted gain
    \begin{equation*}
    p_t = \alpha_t \, \ip{\ghat_C}{\ghat_{\mathrm{upd}}}.
    \end{equation*}
    \item Compute the actual held-out surrogate gain
    \begin{equation*}
    a_t = \widehat{\mathcal L}_C(\theta_{t+1}) - \widehat{\mathcal L}_C(\theta_t).
    \end{equation*}
    \item Form the realization ratio
    \begin{equation*}
    \phit = \frac{a_t}{p_t + \varepsilon}.
    \end{equation*}
    \item Update the slow scale variable
    \begin{equation*}
    c_{t+1} = c_t \exp\!\bigl(\eta_c(\phit - \phi^{\ast})\bigr).
    \end{equation*}
\end{enumerate}
The held-out batch $C$ is required only here.

\section{Why Low-Frequency Calibration is Natural}
Theoretical motivation and engineering motivation point in the same direction:
\begin{itemize}[leftmargin=2em]
    \item $\rhat$ is a fast variable and should respond to the current state of the rollout distribution;
    \item $c_t$ is a slow variable and should not react aggressively to per-step noise.
\end{itemize}
Therefore, there is no reason to update $c_t$ on every step. A very natural implementation is:
\begin{quote}
Estimate $\rhat$ on every step, but update $c_t$ only every $K$ steps.
\end{quote}
When a step is not a calibration step, there is no need to allocate a held-out batch $C$. The sample budget that would otherwise go to $C$ simply does not need to exist on that step.

\section{Implementation Example: $\mathrm{bsz}=256$, $\mathrm{mini\_bsz}=64$, $\mathrm{group\_size}=8$}
Assume the total response budget used in an ordinary training step is 256 responses, with mini-batches of 64 responses and standard GRPO group size 8.

Then the ordinary training batch corresponds to
\begin{equation}
\frac{256}{8} = 32
\end{equation}
prompt groups.

A convenient engineering implementation is:
\begin{itemize}[leftmargin=2em]
    \item split the 32 prompt groups into $A_1$ and $A_2$;
    \item each half contains 16 prompt groups;
    \item each half therefore contains $16 \times 8 = 128$ responses;
    \item with $\mathrm{mini\_bsz}=64$, each half is processed in two mini-batches.
\end{itemize}
So an ordinary step looks like this:
\begin{enumerate}[label=\arabic*., leftmargin=2em]
    \item sample 32 prompt groups under $\theta_t$;
    \item compute $\ghat_{A_1}$ from the first 16 groups;
    \item compute $\ghat_{A_2}$ from the second 16 groups;
    \item form $\rhat$ from the cross-inner-product of the two half-batch gradients;
    \item form $\ghat_{\mathrm{upd}} = (\ghat_{A_1}+\ghat_{A_2})/2$;
    \item set $\alpha_t = c_t\,\rhat$;
    \item update the parameters.
\end{enumerate}
No extra rollout beyond the 256 responses is needed.

On a calibration step, one additionally samples a held-out batch $C$. If $C$ also uses 32 prompt groups with group size 8, then the calibration step incurs one extra block of 256 responses. Since this happens only every $K$ steps, the average rollout multiplier becomes approximately
\begin{equation}
1 + \frac{1}{K}
\end{equation}
relative to the ordinary-step rollout budget used for training itself.

\section{Why This Design Is Better Than a Separate Full ``B'' Stream}
One might imagine using three completely separate streams at every step:
\begin{itemize}[leftmargin=2em]
    \item one for updates,
    \item one for $\rhat$,
    \item one for $\phit$.
\end{itemize}
In GRPO this is unattractive for two reasons:
\begin{enumerate}[label=\arabic*., leftmargin=2em]
    \item if each stream preserves full group size, rollout cost grows rapidly;
    \item if one tries to reduce cost by shrinking each group, advantage quality degrades sharply.
\end{enumerate}
The proposed design avoids both issues:
\begin{itemize}[leftmargin=2em]
    \item $\rhat$ is measured using an \emph{internal split} of the training batch, so there is no dedicated extra rollout just for shape estimation;
    \item only the calibration signal for $c_t$ requires genuinely new held-out data.
\end{itemize}
Thus the only unavoidable measurement overhead is the held-out calibration batch, and even that overhead can be amortized by low-frequency updates.

\section{Practical Safeguards}
Even with the clean role separation above, the ratio $\phit$ may be unstable on some steps. A practical implementation should therefore include safeguards such as:
\begin{enumerate}[label=\arabic*., leftmargin=2em]
    \item update $c_t$ only when $p_t$ is sufficiently positive;
    \item skip the calibration update when $\ip{\ghat_C}{\ghat_{\mathrm{upd}}}$ is non-positive or too small;
    \item clip or smooth $\phit$ before feeding it into the controller;
    \item use a small $\eta_c$ so that $c_t$ remains a slow variable;
    \item optionally impose lower and upper bounds on $c_t$ during early experiments.
\end{enumerate}
These safeguards do not change the logic of the method; they merely keep the controller from chasing batch noise.

\section{Final Recommended Training Loop}
The complete engineering recommendation can be summarized as follows.

\subsection*{Ordinary step}
\begin{enumerate}[label=\arabic*., leftmargin=2em]
    \item Sample one GRPO training batch under the current policy.
    \item Split the batch by prompt groups into two halves.
    \item Compute two half-batch gradients.
    \item Estimate $\rhat$ from their cross-inner-product.
    \item Average the two gradients to form the update gradient.
    \item Use $\alpha_t = c_t\,\rhat$ to update the parameters.
    \item Do not update $c_t$.
\end{enumerate}

\subsection*{Calibration step (every $K$ steps)}
\begin{enumerate}[label=\arabic*., leftmargin=2em]
    \item Perform the ordinary step.
    \item Independently sample a held-out GRPO batch under the old policy.
    \item Compute $p_t$, $a_t$, and $\phit$.
    \item Update $c_t$ using the multiplicative controller.
\end{enumerate}

\subsection*{Core takeaways}
\begin{itemize}[leftmargin=2em]
    \item The samples used to estimate $\rhat$ may and should also contribute to the update gradient.
    \item The only batch that must be genuinely held out is the one used for calibrating $c_t$.
    \item In GRPO, preserving full group integrity is more important than forcing a naive three-way split of responses.
    \item The cleanest first implementation is: \textbf{internal split for $\rhat$, low-frequency held-out calibration for $c_t$}. 
\end{itemize}

\section{Pseudocode}
\begin{algorithm}[h]
\caption{Adaptive learning rate for GRPO: engineering loop}
\begin{algorithmic}[1]
\State Initialize parameters $\theta_0$ and scale $c_0$
\For{$t=0,1,2,\dots$}
    \State Sample one GRPO training batch $A$ under $\theta_t$
    \State Split $A$ by prompt groups into $A_1$ and $A_2$
    \State Compute $\ghat_{A_1}$ and $\ghat_{A_2}$
    \State $\rhat \gets \dfrac{\ip{\ghat_{A_1}}{\ghat_{A_2}}}{(\norm{\ghat_{A_1}}^2 + \norm{\ghat_{A_2}}^2)/2}$
    \State $\ghat_{\mathrm{upd}} \gets (\ghat_{A_1}+\ghat_{A_2})/2$
    \State $\alpha_t \gets c_t\,\rhat$
    \State $\theta_{t+1} \gets \theta_t - \alpha_t\,\ghat_{\mathrm{upd}}$
    \If{$t$ is a calibration step}
        \State Sample held-out batch $C$ under $\theta_t$
        \State Compute $\ghat_C$
        \State $p_t \gets \alpha_t\,\ip{\ghat_C}{\ghat_{\mathrm{upd}}}$
        \State $a_t \gets \widehat{\mathcal L}_C(\theta_{t+1}) - \widehat{\mathcal L}_C(\theta_t)$
        \State $\phit \gets \dfrac{a_t}{p_t + \varepsilon}$
        \State Optionally smooth / clip $\phit$
        \State $c_{t+1} \gets c_t\exp\!\bigl(\eta_c(\phit - 1/2)\bigr)$
    \Else
        \State $c_{t+1} \gets c_t$
    \EndIf
\EndFor
\end{algorithmic}
\end{algorithm}

\end{document}
